\chapter{Conclusion and Limitations}
Atelier Fashion requires consistent SKU identity, immutable purchase evidence, and controlled coordination between stock and money. The investigation links business problems and stakeholder needs to a prioritised requirement baseline. Actor/use-case views distinguish customer goals from internal responsibilities and external provider interfaces. Activity models establish validation and physical inspection, while the order state machine and independent payment/refund states prevent a single status flag from concealing incompatible business facts. Analysis classes distinguish products from variants, carts from reservations, and original order lines from returned quantities.

The proposed design preserves local transaction integrity before external calls and uses durable event/command processing for retries and reconciliation. Exact integer money and deterministic discount allocation support partial refunds. Access and audit controls protect stock and financial decisions, and the complete traceability matrix provides observable release criteria. Native editable Visual Paradigm diagrams and the LaTeX source make the design reviewable and maintainable beyond the PDF.

The study is limited by its fictional business scenario and absence of field elicitation, cost quotations, production workloads, implemented software, and provider acceptance tests. One warehouse, one original shipment, prepaid online payments, and constrained exchanges simplify the first release. A real deployment requires confirmed policies, jurisdictional tax/privacy review, chosen provider contracts, staff training, migration reconciliation, and implementation acceptance testing. Future work may introduce guest checkout, multiple warehouses/shipments, broader exchange pricing, promotions, and richer fit guidance only after revisiting the same consistency and traceability obligations.
